\documentclass[../Main.tex]{subfiles}

\begin{document}
\section{Definition and Examples}
\begin{definition}{Bifurcation}
    Let $\dvec{x} = \vec{f}(\vec{x};\mu)$ be a dynamical system parameterised by $\mu$. A \underline{bifurcation} is a change in the topological structure of the flow as $\mu$ passes through some critical value $\mu_c$. In this case, $\mu_c$ is an \underline{bifurcation point}.
\end{definition}
\begin{remark}
    We will often take $\mu_c = 0$, after possibly shifting it in the definition of $\vec{f}$.
\end{remark}
\begin{example}
    Consider:
    \begin{equation*}
        \dot{x} = \mu x + y - x(x^2 + y^2),\qquad \dot{y} = -x + \mu y - y(x^2 - y^2)
    \end{equation*}
    and transforming this to plane polar coordinates gives:
    \begin{equation*}
        \dot{r} = r(\mu - r^2),\qquad \dot{\theta} = 1.
    \end{equation*}
    We can consider the fixed point at the origin for different values of $\mu$: for $\mu < 0$ we have a stable focus, for $\mu = 0$ we still have a stable focus but slower convergence, and for $\mu > 0$ we have an unstable focus with periodic orbit. Therefore, we conclude that there is a bifurcation at $\mu = 0$.
\end{example}
\begin{remarks}
    \item We will generally be concerned with local bifurcations, looking at topological changes in a small vicinity of fixed points.
    \item We can say that if a fixed point changes stability then there has been a change in topological structure.
        Therefore local bifurcations happen at non-hyperbolic fixed points.
\end{remarks}
\begin{proposition}
    Local bifurcations require the real part of at least one eigenvalue of a fixed point to be zero.
    \label{propBifZeroEVal}
\end{proposition}
\begin{proof}[sketch]
    We can linearise near the fixed point and apply \thmref{propHypMeansLinearOnly} to show that the topological structure changes only when the fixed point becomes non-hyperbolic at the bifurcation point.
\end{proof}
\section{Bifurcations in One Dimension}
\subsection{Saddle-Node and Transcritical Bifurcations}
Consider a 1D system, $\dot{x} = f(x; \mu)$, with a non-hyperbolic fixed point at $x = 0$ when $\mu = 0$. This means $f(0, 0) = 0$ and $\frac{df}{dx}(0, 0) = 0$.

This is related to stationary bifurcations in $\R^n$, where at least one eigenvalue is exactly zero (not just its real part).
\begin{example}
    Consider the dynamical system $\dot{x} = \mu - x^2$. We see that for $\mu < 0$, $\dot{x}$ is always negative, and for $\mu = 0$ we only have the origin where $\dot{x} = 0$. When $\mu > 0$, we will have a region centred around the origin for which $\dot{x} > 0$. See figure~\ref{figOneDBifs} for these cases.
\end{example}
\begin{figure}
    \centering
    \begin{tikzpicture}[scale=1]
        \draw (-2, 2) -- (2, 2)
            node[pos=0, left] {$\mu < 0$}
            pic[pos=0.3, rotate=90] {patharrow}
            pic[pos=0.7, rotate=90] {patharrow};

        \draw (-2, 0) -- (2, 0)
            node[pos=0, left] {$\mu = 0$}
            pic[pos=0.3, rotate=90] {patharrow}
            pic[pos=0.7, rotate=90] {patharrow};
        \fill (0, 0) circle[radius=0.1];

        \draw (-2, -2) -- (2, -2)
            node[pos=0, left] {$\mu < 0$}
            pic[pos=0.2, rotate=90] {patharrow}
            pic[pos=0.5, rotate=-90] {patharrow}
            pic[pos=0.8, rotate=90] {patharrow};
        \fill (-0.5, -2) circle[radius=0.1];
        \fill (0.5, -2) circle[radius=0.1];
    \end{tikzpicture}
    \caption{Diagram of 1D bifurcations}
    \label{figOneDBifs}
\end{figure}

We can expand $f(x;\mu)$ near $x = 0, \mu = 0$:
\begin{align*}
    f(x;\mu) &= f(0;0) + \mu f_\mu(0, 0) + \frac12 f_{xx}(0;0)x^2 + x\mu f_{x\mu}(0;0) + \frac12 f_{\mu\mu}(0;0) + \cdots \\
    &= \left[\mu f_{\mu}(0; 0) + O(\mu^2)\right] + x\mu f_{x\mu}(0; 0) + \frac12 f_{xx}(0;0) + O(x^3) \\
    &= A\mu + B x\mu + Cx^2 + O(\mu^2, x^3)
\end{align*}
where $A = f_\mu(0;0), B = f_{x\mu}(0;0), C = \frac12 f_{xx}(0;0)$. Completing the square:
\begin{equation*}
    \dot{x} = A\mu + C\left(x + \frac{B\mu}{2C}\right)^2 - \frac{B^2 \mu^2}{4C} + O(\mu^2, x^3)
\end{equation*}
but we can remove this last term because we are working to $O(\mu)$:
\begin{equation*}
    \dot{x} \approx A\mu + C\left(x + \frac{B\mu}{2C}\right)^2.
\end{equation*}
Transforming coordinates as:
\begin{equation*}
    x \mapsto -y + \frac{B}{2C}\mu,\qquad t \mapsto T / C, \qquad \mu \mapsto \tilde{\mu} = -\frac{A}{C}\mu
\end{equation*}
yields the \underline{normal form}:
\begin{equation*}
    \dot{y} = \tilde\mu - y^2.
\end{equation*}
This requires $A$ and $C$ to be non-zero. This type of bifurcation is known as a \underline{saddle-node} bifurcation.

If in fact $A = 0$ we have a \underline{transcritical bifurcation}, with normal form given by:
\begin{equation*}
    \dot{y} = \tilde{\mu} y - y^2
\end{equation*}
under the transformations $y = -x, T = Ct, \tilde\mu = \frac B C \mu$. This gives a stable fixed point at the origin (and an unstable one at $x = \mu$) for $\mu < 0$, a half-stable, half-unstable fixed point at the origin for $\mu = 0$, and an unstable fixed point at the origin (and stable one at $x = \mu$) when $\mu > 0$. We can represent this as a plot of $\mu$ against $x$ -- see figure~\ref{figTranscritBif}. Here the dotted line represents an unstable fixed point, the solid line represents a stable one.
\begin{figure}
    \centering
    \begin{tikzpicture}[scale=1]
        \draw[->] (-2, 0) -- (2, 0) node[right] {$\mu$};
        \draw (0, -2) -- (0, 2) node[above] {$x$};

        \draw[dashed, blue] (-2, -2) -- (0, 0);
        \draw[blue] (-2, 0) -- (0, 0);

        \draw[dashed, blue] (0, 0) -- (2, 0);
        \draw[blue] (0, 0) -- (2, 2);
    \end{tikzpicture}
    \caption{Diagram for a transcritical bifurcation}
    \label{figTranscritBif}
\end{figure}

\subsection{Pitchfork Bifurcations}
If the dynamical system has a symmetry, and is unchanged under $x \mapsto -x$, then our previous expansion for $\dot{x}$ requires $A = C = 0$. To see the dynamics of the bifurcation we go to higher order in $x$ and $\mu$,
\begin{equation*}
    \dot{x} = \mu x f_{x\mu}(0;0) + \frac16 x^3 f_{xxx}(0;0) + \cdots
\end{equation*}
which, after a change of variables, has normal form:
\begin{equation*}
    \dot{x} = \mu x - ax^3
\end{equation*}
for $a = \pm 1$. In the case of $a = 1$, we have a \underline{supercritical pitchfork bifurcation}. This has fixed points at $x = \pm \sqrt{\mu}$ as well as $x = 0$.

In the case of $a = -1$, we have a \underline{subcritical pitchfork bifurcation}, with fixed points at $\pm \sqrt{-\mu}$. For both types of pitchfork, see figure~\ref{figPitchforkBifs}.
\begin{figure}
    \centering
    \begin{tikzpicture}[scale=1]
        \node {TODO}; %TODO
    \end{tikzpicture}
    \caption{Diagram of supercitical and subcritical pitchfork bifurcations}
    \label{figPitchforkBifs}
\end{figure}
\subsection{Structural Stability of One-Dimensional Bifurcations}
We will consider perturbed dynamical systems of the form:
\begin{equation*}
    \dot{x} = f(x;\mu) + \epsilon(x)
\end{equation*}
where we can assume that $f$ is in normal form.

Consider a saddle node, with $\dot{x} = \mu - x^2$.
\begin{itemize}
    \item If we add $\epsilon$ to the right-hand side,
        \begin{equation*}
            \dot{x} = \epsilon + \mu - x^2
        \end{equation*}
        which we can see does not change the dynamics (if we redefine $\mu$).
    \item If we add $\epsilon x$ to the right-hand side,
        \begin{equation*}
            \dot{x} = \mu + \epsilon x - x^2
        \end{equation*}
        and we can complete the square, rescale $x$ and redefine $\mu$ to get back to the same normal form.
\end{itemize}
Since our expansion worked to order $O(x)$, we conclude that the saddle-node is structurally stable.

Consider instead a transcritical bifurcation, with normal form $\dot{x} = \mu x - x^2$. If we perturb this for small $\epsilon$, we get $\dot{x} = \epsilon + \mu x - x^2$.
\begin{itemize}
    \item If $\epsilon < 0$, we find a region of $\mu$ around $0$ where there are no fixed points.
    \item If $\epsilon > 0$, we find that the fixed points never change stability and there is no bifurcation at all!
\end{itemize}

Now if we have a pitchfork bifurcation, with normal form $\dot{x} = \mu x - x^3$, we also find it is structurally unstable under adding a constant $\epsilon$.
\section{Bifurcations in Two Dimensions}
Consider the generalisation of what we were doing before,
\begin{equation*}
    \dvec{x} = \vec{f}(\vec{x};\mu)
\end{equation*}
We can have four possibilities depending on the form of the linearised matrix $A$ around the origin.
\begin{itemize}
    \item One eigenvalue is zero -- a \underline{simple bifurcation} or \underline{stationary bifurcation} in 2D, and it is essentially the same as our one-dimensional bifurcations.
    \item The eigenvalues are imaginary -- $\lambda = \pm i\omega$ -- an \underline{oscillatory bifurcation} or \underline{Hopf bifurcation}.
    \item Both eigenvalues are zero and non-degenerate -- a \underline{double-zero}\\\underline{bifurcation}.
    \item Both eigenvalues are zero and degenerate -- a \underline{Takens-Bogdanov} bifurcation.
\end{itemize}
The last two cases will not be covered in this course.
\subsection{The Centre Manifold}
Consider a system:
\begin{equation}
    \dot{x} = xy - \mu x,\qquad \dot{y} = -y + x^2
    \label{eqnBifSystem}
\end{equation}
There is a fixed point at the origin for any value of $\mu$, and for $\mu > 0$ we find a further two fixed points at $(\pm \sqrt{\mu}, \mu)$. Linearising around the origin gives:
\begin{equation*}
    A = 
    \begin{pmatrix}
        -\mu & 0 \\
        0 & -1
    \end{pmatrix}
\end{equation*}
and so when $\mu = 0$, the origin is a non-hyperbolic fixed point, so the dynamical system admits a bifurcation. We analyse close to the origin and bifurcation point by requiring $\abs x \sim \abs y \sim \abs \mu \ll 1$. Then $\abs{\dot x} \sim O(\mu^2)$ and $\abs{\dot y} \sim O(\mu)$, so the system evolves much faster in $y$ than $x$ towards $y = x^2$. On this curve, we find:
\begin{equation*}
    \dot{x} = x(x^2 - \mu)
\end{equation*}
and we have reduced our two-dimension problem to a 1-dimensional one. We recognise this as a subcritical pitchfork.
\end{document}